% Pista ana-23s-q4 · Anàlisi
% Procedència: PAU Catalunya 2023 · Convocatòria de setembre · Sèrie 2 · Qüestió 4
\documentclass[11pt,a4paper]{article}
\usepackage{pista}
\pistainfo{Catalunya 2023}{Convocatòria de setembre}{Sèrie 2}

\begin{document}

\section*{\textcolor{blauPista}{Qüestió 4}}

\noindent Sigui la funció $f(x)$ definida per $f(x) = -3x + e^{2x^{3} - 1}$.

\subsection*{\textit{a)} \normalfont Justifiqueu que $f(x) = 2$ té una solució en l'interval $(-1, 0)$. \hspace*{\fill}\textnormal{[1,25~punts]}}

\pista{Reformula la pregunta: vols saber si s'anul·la la funció auxiliar
\[
g(x) = f(x) - 2
\]
En primer lloc, justifica que $g$ és contínua a $[-1, 0]$: és la suma d'una funció polinòmica i d'una exponencial, totes dues contínues. A continuació, calcula $g(-1)$ i $g(0)$. Si tenen signes oposats, el teorema de Bolzano assegura que $g$ s'anul·la en algun valor de l'interval $(-1, 0)$, és a dir, que en aquest valor $f(x) = 2$.}

\subsection*{\textit{b)} \normalfont Sigui la funció $h(x) = -3x^{2} + e^{2x^{3} - 1}$. Calculeu l'àrea de la regió compresa entre les gràfiques de les funcions $f(x)$ i $h(x)$. \hspace*{\fill}\textnormal{[1,25~punts]}}

\pista{En primer lloc, calcula la diferència $f(x) - h(x)$ i simplifica-la. No és tan difícil com sembla, perquè les dues exponencials s'anul·len.

\medskip
A continuació, resol $f(x) - h(x) = 0$: les solucions et marquen precisament l'interval sobre el qual integres. Finalment, calcula l'àrea amb el valor absolut de la integral:
\[
\text{Àrea} = \left| \int_{x_{1}}^{x_{2}} \bigl(f(x) - h(x)\bigr)\,dx \right|
\]}

\end{document}
